% UIB Advanced Algorithms Report Template
% Based on IEEE conference format with custom modifications for course reports
% This template includes the formatting customizations from P1

\documentclass[conference,compsoc]{IEEEtran}

% ===== PACKAGES =====
\usepackage{amsmath,amsfonts,amssymb}
\usepackage{algorithmic}
\usepackage{algorithm}
\usepackage{array}
\usepackage[caption=false,font=footnotesize,labelfont=sf,textfont=sf]{subfig}
\usepackage{textcomp}
\usepackage{stfloats}
\usepackage{url}
\usepackage{verbatim}
\usepackage{graphicx}
\usepackage[nocompress]{cite}
\usepackage{listings}
\usepackage{xcolor}
\usepackage[catalan]{babel}
\usepackage[T1]{fontenc}
\usepackage[utf8]{inputenc}
\usepackage{booktabs}
\usepackage{float}
\usepackage{cuted}

% ===== CATALAN LANGUAGE CUSTOMIZATION =====
% Rename abstract and keywords to Catalan
\renewcommand{\abstractname}{Resum}
\renewcommand{\IEEEkeywordsname}{Termes d'índex}

% ===== ABSTRACT AND KEYWORDS FORMATTING =====
% Custom formatting: bold italic labels, normal-weight body text
\makeatletter
\def\abstract{\normalfont\@IEEEtweakunitybaselinestretch{1.15}%
    \if@twocolumn
      \@IEEEabskeysecsize\noindent\textbf{\textit{\abstractname}}---\normalfont\@IEEEabskeysecsize\relax
    \else
      \bgroup\par\addvspace{0.5\baselineskip}\centering\vspace{-1.78ex}\@IEEEabskeysecsize\textbf{\abstractname}\par\addvspace{0.5\baselineskip}\egroup\quotation\normalfont\@IEEEabskeysecsize%
    \fi\@IEEEgobbleleadPARNLSP}
\def\IEEEkeywords{\normalfont\@IEEEtweakunitybaselinestretch{1.15}%
    \if@twocolumn
      \@IEEEabskeysecsize\vskip 0.5\baselineskip plus 0.25\baselineskip minus 0.25\baselineskip\noindent
      \textbf{\textit{\IEEEkeywordsname}}---\normalfont\@IEEEabskeysecsize\relax
    \else
      \bgroup\par\addvspace{0.5\baselineskip}\centering\@IEEEabskeysecsize\textbf{\IEEEkeywordsname}\par\addvspace{0.5\baselineskip}\egroup\quotation\normalfont\@IEEEabskeysecsize%
    \fi\@IEEEgobbleleadPARNLSP}
\makeatother

% ===== CODE LISTING CONFIGURATION =====
% Configuration for Java code listings (adjust language as needed)
\lstset{
    language=Java,
    basicstyle=\ttfamily\footnotesize,
    keywordstyle=\color{blue}\bfseries,
    commentstyle=\color{gray}\itshape,
    stringstyle=\color{red!70!black},
    numbers=left,
    numberstyle=\tiny\color{gray},
    numbersep=5pt,
    frame=single,
    breaklines=true,
    showstringspaces=false,
    tabsize=4,
    captionpos=b
}

% ===== HYPHENATION RULES =====
% Add custom hyphenation rules for Catalan technical terms as needed
\hyphenation{a-sin-to-tic a-sin-to-tics al-go-ris-me al-go-ris-més con-tro-la-dor mul-ti-pli-ca-ti-va}

% ===== DOCUMENT BEGINS =====
\begin{document}

% ===== TITLE AND AUTHOR =====
\title{[Títol de la Pràctica] \\ Pràctica X}

\author{\IEEEauthorblockN{[Nom de l'Estudiant]}
\IEEEauthorblockA{Universitat de les Illes Balears (UIB)\\
Grau en Enginyeria Informàtica\\
Algorismes Avançats, Curs 2025--2026}}

\maketitle

% ===== ABSTRACT AND KEYWORDS =====
% Using strip environment for full-width abstract in conference format
\begin{strip}
\begin{abstract}
[Escriu aquí el resum de la pràctica. Ha de ser un paràgraf concís que
descrigui els objectius, la metodologia i els resultats principals.]
\end{abstract}

\begin{IEEEkeywords}
[paraula clau 1], [paraula clau 2], [paraula clau 3], [paraula clau 4]
\end{IEEEkeywords}
\end{strip}

\IEEEpeerreviewmaketitle

% ===== INTRODUCTION =====
\section{Introducció}
[Introdueix el context de la pràctica, els objectius i l'estructura
de la memòria.]

% ===== THEORETICAL FOUNDATIONS =====
\section{Fonaments Teòrics}

\subsection{[Subtítol 1]}
[Contingut teòric relacionat amb la pràctica.]

\subsection{[Subtítol 2]}
[Més contingut teòric.]

% Example of equation
\begin{equation}
\label{eq:example}
T(n) = c \cdot f(n)
\end{equation}

% Example of table
\begin{table}[!h]
      \caption{Exemple de taula.\label{tab:example}}
      \centering
      \begin{tabular}{r|rrr}
            \toprule
            $n$ & $f_1(n)$ & $f_2(n)$ & $f_3(n)$ \\
            \midrule
            10     & 10       & 100         & 1\,000 \\
            100    & 100      & 10\,000     & $10^6$ \\
            1\,000  & 1\,000  & $10^6$      & $10^9$ \\
            \bottomrule
      \end{tabular}
\end{table}

% ===== DESIGN AND IMPLEMENTATION =====
\section{Disseny i Implementació}

\subsection{[Component 1]}
[Descripció del component.]

% Example of code listing
\begin{lstlisting}[caption={Exemple de codi.}]
public class Example {
    public static void main(String[] args) {
        System.out.println("Hello, World!");
    }
}
\end{lstlisting}

\subsection{[Component 2]}
[Més descripció.]

% Example of figure
\begin{figure}[!h]
    \centering
    \includegraphics[width=\columnwidth]{example-figure.png}
    \caption{Exemple de figura.\label{fig:example}}
\end{figure}

% For full-width figure in conference format, use figure* environment
% \begin{figure*}[!h]
%     \centering
%     \includegraphics[width=\textwidth]{wide-figure.png}
%     \caption{Exemple de figura ampla.\label{fig:wide}}
% \end{figure*}

% ===== RESULTS AND DISCUSSION =====
\section{Resultats i Discussió}

[Presenta els resultats experimentals i discuteix-los. Referencia les
figures i taules adequadament, per exemple: la Fig.~\ref{fig:example}
mostra que... o segons la Taula~\ref{tab:example}...]

% Example of fixed-position figure
\begin{figure}[H]
    \centering
    \includegraphics[width=\columnwidth]{results-figure.png}
    \caption{Resultats experimentals.\label{fig:results}}
\end{figure}

% ===== CONCLUSIONS =====
\section{Conclusions}

[Resume els resultats principals, avalua l'assoliment dels objectius
i proposa treball futur.]

% ===== BIBLIOGRAPHY =====
\begin{thebibliography}{9}
\bibliographystyle{IEEEtran}

\bibitem{ref_example1}
A. Author, B. Coauthor, and C. Thirdauthor, \textit{Title of the Book}.
City, State, Country: Publisher, Year.

\bibitem{ref_example2}
D. Researcher, ``Title of the Article,'' \textit{Journal Name}, vol.~X,
no.~Y, pp.~123--456, Month Year.

\bibitem{ref_example3}
E. WebAuthor, ``Title of Web Resource,'' Organization, Year. [Online].
Available: \url{https://example.com/resource}

% Common references for Advanced Algorithms course:
\bibitem{ref_cormen}
T. H. Cormen, C. E. Leiserson, R. L. Rivest, and C. Stein,
\textit{Introduction to Algorithms}, 4th~ed. Cambridge, MA, USA:
MIT Press, 2022.

\bibitem{ref_sedgewick}
R. Sedgewick and K. Wayne, \textit{Algorithms}, 4th~ed. Upper Saddle
River, NJ, USA: Addison-Wesley, 2011.

\bibitem{ref_knuth}
D. E. Knuth, \textit{The Art of Computer Programming, Vol. 1:
Fundamental Algorithms}, 3rd~ed. Reading, MA, USA: Addison-Wesley, 1997.

\bibitem{ref_mvc}
E. Gamma, R. Helm, R. Johnson, and J. Vlissides, \textit{Design Patterns:
Elements of Reusable Object-Oriented Software}. Reading, MA, USA:
Addison-Wesley, 1994.

\bibitem{ref_observer}
F. Buschmann, R. Meunier, H. Rohnert, P. Sommerlad, and M. Stal,
\textit{Pattern-Oriented Software Architecture: A System of Patterns}.
Chichester, U.K.: Wiley, 1996.

\bibitem{ref_swing}
J. Elliott, R. Eckstein, M. Loy, D. Wood, and B. Cole, \textit{Java
Swing}, 2nd~ed. Sebastopol, CA, USA: O'Reilly Media, 2002.

\bibitem{ref_java}
Oracle Corporation, ``The Java Tutorials --- Concurrency,''
[Online]. Available:
\url{https://docs.oracle.com/javase/tutorial/essential/concurrency/}

\end{thebibliography}

\end{document}
